% ========================================================
% BÀI TẬP TOÁN 5 - KẾT NỐI TRI THỨC VỚI CUỘC SỐNG
% Bài 57: Cộng, trừ số đo thời gian
% Định dạng: BTTeX Part III - Trắc nghiệm điền số ngắn
% ========================================================

% Câu 1
\begin{ex}
  Tính: $3\text{ giờ } 15\text{ phút} + 2\text{ giờ } 30\text{ phút}$. Kết quả có số phút là bao nhiêu?
  \shortans{45}
  \loigiai{
    15 + 30 = 45 phút.
  }
\end{ex}

% Câu 2
\begin{ex}
  Tính: $4\text{ giờ } 35\text{ phút} + 1\text{ giờ } 40\text{ phút}$. Kết quả đổi thành bao nhiêu giờ $15$ phút?
  \shortans{6}
  \loigiai{
    $5\text{ giờ } 75\text{ phút} = 6\text{ giờ } 15\text{ phút}$.
  }
\end{ex}

% Câu 3
\begin{ex}
  Tính: $2\text{ ngày } 14\text{ giờ} + 3\text{ ngày } 18\text{ giờ}$. Kết quả có số giờ dư (sau khi đổi ngày) là bao nhiêu?
  \shortans{8}
  \loigiai{
    $14 + 18 = 32\text{ giờ} = 1\text{ ngày } 8\text{ giờ}$. Số giờ là 8.
  }
\end{ex}

% Câu 4
\begin{ex}
  Tính: $5\text{ phút } 20\text{ giây} + 3\text{ phút } 45\text{ giây}$. Kết quả đổi thành bao nhiêu phút $5$ giây?
  \shortans{9}
  \loigiai{
    $8\text{ phút } 65\text{ giây} = 9\text{ phút } 5\text{ giây}$.
  }
\end{ex}

% Câu 5
\begin{ex}
  Tính: $7\text{ giờ } 45\text{ phút} - 3\text{ giờ } 20\text{ phút}$. Kết quả có số phút là bao nhiêu?
  \shortans{25}
  \loigiai{
    45 - 20 = 25 phút.
  }
\end{ex}

% Câu 6
\begin{ex}
  Tính: $5\text{ giờ } 15\text{ phút} - 2\text{ giờ } 40\text{ phút}$. Đổi $5\text{ giờ } 15\text{ phút} = 4\text{ giờ } 75\text{ phút}$. Kết quả có số phút là bao nhiêu?
  \shortans{35}
  \loigiai{
    75 - 40 = 35 phút.
  }
\end{ex}

% Câu 7
\begin{ex}
  Tính: $4\text{ phút } 10\text{ giây} - 1\text{ phút } 35\text{ giây}$. Kết quả có số giây là bao nhiêu?
  \shortans{35}
  \loigiai{
    $3\text{ phút } 70\text{ giây} - 1\text{ phút } 35\text{ giây} = 2\text{ phút } 35\text{ giây}$.
  }
\end{ex}

% Câu 8
\begin{ex}
  Một người đi từ $A$ lúc $7\text{ giờ } 15\text{ phút}$ và đến $B$ lúc $9\text{ giờ } 45\text{ phút}$. Thời gian đi hết bao nhiêu phút?
  \shortans{150}
  \loigiai{
    $9\text{ giờ } 45\text{ phút} - 7\text{ giờ } 15\text{ phút} = 2\text{ giờ } 30\text{ phút} = 150$ phút.
  }
\end{ex}

% Câu 9
\begin{ex}
  Một trận bóng đá bắt đầu lúc $15\text{ giờ } 30\text{ phút}$, kéo dài $90$ phút và nghỉ giữa hiệp $15$ phút. Trận đấu kết thúc lúc mấy giờ $15$ phút?
  \shortans{17}
  \loigiai{
    Tổng thời gian là $105$ phút = $1\text{ giờ } 45\text{ phút}$. Kết thúc lúc $15\text{h}30 + 1\text{h}45 = 17\text{ giờ } 15\text{ phút}$.
  }
\end{ex}

% Câu 10
\begin{ex}
  Một ô tô khởi hành từ $A$ lúc $6\text{ giờ } 30\text{ phút}$, đi hết $2\text{ giờ } 45\text{ phút}$. Xe đến $B$ lúc mấy giờ $15$ phút?
  \shortans{9}
  \loigiai{
    $6\text{h}30 + 2\text{h}45 = 8\text{h}75 = 9\text{ giờ } 15\text{ phút}$.
  }
\end{ex}

% Câu 11
\begin{ex}
  Tính: $12\text{ năm } 6\text{ tháng} - 4\text{ năm } 8\text{ tháng}$. Kết quả có số tháng là bao nhiêu?
  \shortans{10}
  \loigiai{
    $11\text{ năm } 18\text{ tháng} - 4\text{ năm } 8\text{ tháng} = 7\text{ năm } 10\text{ tháng}$.
  }
\end{ex}

% Câu 12
\begin{ex}
  Tính: $8\text{ ngày } 6\text{ giờ} - 3\text{ ngày } 12\text{ giờ}$. Kết quả có số giờ là bao nhiêu?
  \shortans{18}
  \loigiai{
    $7\text{ ngày } 30\text{ giờ} - 3\text{ ngày } 12\text{ giờ} = 4\text{ ngày } 18\text{ giờ}$.
  }
\end{ex}

% Câu 13
\begin{ex}
  Bạn Lan học bài từ $19\text{ giờ } 30\text{ phút}$ đến $21\text{ giờ } 10\text{ phút}$. Lan học trong bao nhiêu phút?
  \shortans{100}
  \loigiai{
    $21\text{h}10 - 19\text{h}30 = 1\text{ giờ } 40\text{ phút} = 100$ phút.
  }
\end{ex}

% Câu 14
\begin{ex}
  Một người thợ làm sản phẩm thứ nhất hết $1\text{ giờ } 20\text{ phút}$, làm sản phẩm thứ hai hết $1\text{ giờ } 35\text{ phút}$. Làm cả hai sản phẩm hết bao nhiêu phút?
  \shortans{175}
  \loigiai{
    $1\text{h}20 + 1\text{h}35 = 2\text{h}55 = 175$ phút.
  }
\end{ex}

% Câu 15
\begin{ex}
  Ở câu trên, sản phẩm thứ hai làm lâu hơn sản phẩm thứ nhất bao nhiêu phút?
  \shortans{15}
  \loigiai{
    $35 - 20 = 15$ phút.
  }
\end{ex}

% Câu 16
\begin{ex}
  Một người đi xe máy từ quê ra phố hết $3\text{ giờ } 15\text{ phút}$. Khi về người đó đi ô tô chỉ hết $2\text{ giờ } 30\text{ phút}$. Thời gian về nhanh hơn thời gian đi bao nhiêu phút?
  \shortans{45}
  \loigiai{
    $3\text{h}15 - 2\text{h}30 = 45$ phút.
  }
\end{ex}

% Câu 17
\begin{ex}
  Một cuộc thi chạy tiếp sức có 3 vận động viên: VĐV 1 chạy hết $45$ giây, VĐV 2 hết $48$ giây, VĐV 3 hết $47$ giây. Cả đội chạy hết bao nhiêu phút?
  \shortans{140}
  \loigiai{
    $45 + 48 + 47 = 140$ giây.
  }
\end{ex}

% Câu 18
\begin{ex}
  Tính: $15\text{ giờ } 20\text{ phút} + 4\text{ giờ } 50\text{ phút}$. Kết quả đổi thành bao nhiêu giờ $10$ phút?
  \shortans{20}
  \loigiai{
    $19\text{h}70 = 20\text{ giờ } 10\text{ phút}$.
  }
\end{ex}

% Câu 19
\begin{ex}
  Một tàu hỏa rời ga Hà Nội lúc $22\text{ giờ}$ đêm hôm trước và đến Vinh lúc $4\text{ giờ } 30\text{ phút}$ sáng hôm sau. Thời gian tàu chạy hết bao nhiêu phút?
  \shortans{390}
  \loigiai{
    Từ 22h đến 24h là 2 giờ; từ 0h đến 4h30 là 4h30. Tổng là $6\text{ giờ } 30\text{ phút} = 390$ phút.
  }
\end{ex}

% Câu 20
\begin{ex}
  Tính: $10\text{ phút } - 3\text{ phút } 40\text{ giây}$. Kết quả có số giây là bao nhiêu?
  \shortans{20}
  \loigiai{
    $9\text{ phút } 60\text{ giây} - 3\text{ phút } 40\text{ giây} = 6\text{ phút } 20\text{ giây}$.
  }
\end{ex}

% Câu 21
\begin{ex}
  Tính: $6\text{ giờ } - 2\text{ giờ } 15\text{ phút}$. Kết quả có số phút là bao nhiêu?
  \shortans{45}
  \loigiai{
    $5\text{ giờ } 60\text{ phút} - 2\text{ giờ } 15\text{ phút} = 3\text{ giờ } 45\text{ phút}$.
  }
\end{ex}

% Câu 22
\begin{ex}
  Một ca làm việc bắt đầu lúc $8\text{ giờ}$ và kết thúc lúc $16\text{ giờ } 30\text{ phút}$, trong đó nghỉ trưa $1\text{ giờ } 30\text{ phút}$. Thời gian làm việc thực tế là bao nhiêu giờ?
  \shortans{7}
  \loigiai{
    Tổng thời gian là 8h30. Trừ nghỉ 1h30 còn 7 giờ.
  }
\end{ex}

% Câu 23
\begin{ex}
  Một chương trình truyền hình dự kiến chiếu $1\text{ giờ } 15\text{ phút}$ nhưng bị lố $18$ phút. Chương trình chiếu hết bao nhiêu phút?
  \shortans{93}
  \loigiai{
    $75 + 18 = 93$ phút.
  }
\end{ex}

% Câu 24
\begin{ex}
  Tính: $3\text{ ngày } 8\text{ giờ} + 1\text{ ngày } 16\text{ giờ}$. Kết quả là đúng bao nhiêu ngày?
  \shortans{5}
  \loigiai{
    $4\text{ ngày } 24\text{ giờ} = 5\text{ ngày}$.
  }
\end{ex}

% Câu 25
\begin{ex}
  Tính: $7\text{ giờ } 12\text{ phút} - 2\text{ giờ } 30\text{ phút}$. Kết quả có số phút là bao nhiêu?
  \shortans{42}
  \loigiai{
    $6\text{h}72 - 2\text{h}30 = 4\text{h}42$. Số phút là 42.
  }
\end{ex}
